\chapter{Errores y códigos}
\label{sec:errores}

\section{HTTP API}

El daemon sigue códigos HTTP estándar. Cada error trae un body JSON:

\begin{lstlisting}[basicstyle=\ttfamily\small]
{
  "error":   "dsl_parse",
  "message": "unexpected token at line 1, col 14",
  "details": { ... }
}
\end{lstlisting}

\begin{table}[H]
\centering
\small
\begin{tabular}{@{}lll@{}}
\toprule
Código & \code{error} & Causa típica \\
\midrule
400 & \code{bad_request}        & DTO malformado, tipo incorrecto. \\
401 & \code{unauthorized}       & Header \code{Authorization} ausente o token inválido. \\
403 & \code{forbidden}          & Feature flag desactivado. \\
404 & \code{not_found}          & Nodo, hunt, nota o sesión inexistente. \\
409 & \code{conflict}           & Ingesta en curso, sesión bloqueada. \\
422 & \code{dsl_parse}          & DSL inválido. \\
422 & \code{schema_drift}       & Batch viola tolerancia de drift. \\
422 & \code{invariant_failed}   & Grafo hipotético rompe invariantes. \\
500 & \code{internal}           & Error no clasificado. Revisar log del daemon. \\
\bottomrule
\end{tabular}
\caption{Códigos HTTP emitidos por la API.}
\end{table}

\section{MCP tools}

Cada tool MCP propaga el error de la API subyacente como excepción
tipada. El cliente recibe:

\begin{lstlisting}[basicstyle=\ttfamily\small]
{
  "isError": true,
  "content": [{
    "type": "text",
    "text": "MCP error: unauthorized (401) — GRAPHHUNTER_API_TOKEN missing or invalid"
  }]
}
\end{lstlisting}

Errores específicos de tools (antes de llegar al API):

\begin{itemize}
    \item \textbf{Zod validation}: input no cumple \code{inputSchema}
    de la tool. Mensaje incluye el path Zod fallido.
    \item \textbf{Timeout}: 30s default (lectura), 300s (ingest /
    hunt). Tool devuelve \code{timeout}.
    \item \textbf{Connection refused}: daemon no corre o URL mal
    configurada. Mensaje incluye \code{GRAPHHUNTER_API_URL}
    efectivo.
    \item \textbf{HTML en lugar de JSON}: otro servicio ocupa el
    puerto. Error \code{not_graphhunter}.
\end{itemize}

\section{DSL parser}

Errores emitidos por \code{POST /parse_dsl} y por
\code{POST /run_hunt}:

\begin{table}[H]
\centering
\small
\begin{tabular}{@{}lp{7cm}@{}}
\toprule
\code{error} & Significado \\
\midrule
\code{unexpected_token}     & Token no esperado en esa posición. \\
\code{unknown_entity_type}  & Tipo no registrado en el
\code{ExtensionRegistry}. \\
\code{unknown_relation}     & Ídem, relación. \\
\code{unknown_predicate}    & Predicado no implementado. \\
\code{invalid_duration}     & Formato de \code{window} incorrecto. \\
\code{empty_hypothesis}     & La hipótesis no tiene steps. \\
\bottomrule
\end{tabular}
\end{table}

Todos los errores DSL traen \code{line} y \code{column}.

\section{C ABI}

La ABI de \filepath{libgraphmatch} no usa códigos de error; usa
contratos (ver cap.~\ref{sec:cabi}):

\begin{itemize}
    \item \code{gm_matcher_run(m, g)}: devuelve \code{NULL} si
    \code{m} mal configurado o \code{g} no finalizado.
    \item \code{gm_results_get_*}: índice fuera de rango devuelve
    \code{NULL} y \code{*out_len = 0}.
    \item Mismatch de \code{gm_abi_version}: el cliente Rust aborta
    con panic (el binario no arranca).
    \item Violación de invariantes (doble free, uso después de
    free): UB. El wrapper Rust \code{GmGraph} previene esto por
    typestate y RAII.
\end{itemize}

\section{Ingesta}

La ingesta tiene errores granulares (ver también
Dead Letter Queue, \S\ref{sec:api-http}):

\begin{table}[H]
\centering
\small
\begin{tabular}{@{}lp{7cm}@{}}
\toprule
\code{error} & Significado \\
\midrule
\code{parser_not_found}    & Ningún parser reconoce la muestra. \\
\code{mapping_missing_time} & \code{FieldConfig} sin rol
\code{TIME} obligatorio. \\
\code{time_parse_failed}   & Ningún formato de timestamp matchea. \\
\code{dlq_enqueued}        & Evento parseable pero violaba invariante;
guardado en DLQ. \\
\code{session_locked}      & Otra ingesta está corriendo en esa
sesión. \\
\bottomrule
\end{tabular}
\end{table}

Los eventos en DLQ se listan con \code{POST /list_dead_letters} y
se reintentan con \code{POST /reingest_dead_letter}.

\section{Convención de log}

Todos los errores relevantes se loguean en el daemon con nivel
\code{error} o \code{warn}. Correlación vía \code{request_id}
inyectado en cada request HTTP (header
\code{X-Request-Id}). Revisar \filepath{app.log} para detalle.
